\section{Introduction}
\label{sec:intro}

\begin{figure}[ht]
    \centering
    \includegraphics[width=0.95\textwidth]{Figures/working-memory-illustration.png}
    \caption{Conceptual illustration of information storage in working memory. The neural network can store and recall features of an item. The volatile nature of the synapses allows the memory of the stored features to fade in time. After depletion of the memory, the features of a different item can be stored without significant interferences.}
    \label{fig:fig1}
\end{figure}

The short-term storage of information, one of the main properties of \gls{wm}, is at the forefront of most cognitive abilities in humans and animals \cite{barak2014working, mi2017synaptic}. Tasks that involve \gls{wm} include visual processing, speech comprehension, and episodic planning; thus impairment of \gls{wm} results in a partial or complete loss of these abilities \cite{cowan2010magical,miller2018working}. Several models have been suggested to explain the physiological mechanisms that underlie the \gls{wm} in the brain. Today it is becoming increasingly evident that a mixture of specialized mechanisms enable an ensemble of versatile memory systems in the brain that cover memory duration from several seconds to minutes \cite{mongillo2008synaptic, miller2018working}.

One such mechanism for \gls{wm} is based on short-term dynamics of synapses, i.e., \gls{stp}. Figure~\ref{fig:fig1} illustrates the synaptic model of \gls{wm}. When a new memory item is stored in the \gls{wm} network, an ensemble of neurons that represents the item is activated and synaptic strengths between these neurons are transiently potentiated. This transient increase in synaptic strength is caused by short term mechanisms that put the synapse in a state of higher effectiveness \cite{mi2017synaptic, mongillo2008synaptic}. This state of high effectiveness supports the system to retrieve the stored memory item. %The stored memory item can be retrieved during a brief time window in which the synapses are still in the state of high effectiveness. 
The transient decay of the synaptic strengths and, thus, of the memory items makes the \gls{wm} susceptible for new content.

Various mechanisms for \gls{wm} were introduced also in modern \gls{ml} to improve the performance in complex sequence processing tasks including natural language \cite{weston2014memory, graves2014neural, vinyals2015pointer, santoro2016meta}. Indeed, similar to the \gls{wm}, \glspl{rnn} provide the possibility to perform associative memory and recall stored information thanks to the presence of internal feedback loops that ensures the persistence of data. In recent studies, recurrency demonstrated to be a crucial aspect in, e.g., grammar learning \cite{alamia2020SciRep}. In contrast to feedforward architectures that are often being used in \gls{ml}, \glspl{rnn} inhere the desirable computational properties to represent long-term dependencies in sequential data, due to their computational load \cite{dehghani2018universal, schafer2006recurrent}.

Network architectures such as \gls{lstm} and \glspl{gru} try to optimize the use of hardware resources while still exploiting recurrency \cite{yu2019NeurComp}. However, these solutions are still expensive in terms of hardware area and number of operations. Bespoke hardware optimizing resources in terms of power consumption, area, and computational workload could therefore facilitate the implementation of \glspl{rnn}. At present, \gls{cmos}-only solutions include digital standard hardware, e.g., \glspl{gpu} and \glspl{fpga}, as well as analog or mixed-signal \gls{asic}. Digital circuits usually imply a heavy computational workload and the need of waiting for several clock cycles \cite{amirsoleimani2020AdvIntellSys}, while \glspl{asic} have a higher energy efficient approach by adopting subthreshold circuits working asynchronously \cite{rahimi2020AdvIntellSys}. However, the temporal dynamics are usually implemented by charging or discharging a capacitor with a constant current \cite{chicca2014ProcIEEE}, thus consuming power. Furthermore, when dealing with biologically relevant time scales, the size of the capacitors becomes non-negligible.

Alternative approaches with better scaling and lower energy consumption are deeply sought to enable neuromorphic circuits with bio-realistic \glspl{wm} \cite{covi2021FrontNeurosci}. In this respect, a promising technology is represented by memristive devices \cite{christensen2022NCE}, namely two-terminal devices able to reversibly change their conductance upon the application of proper electrical stimuli. Memristive devices show a broad range of attractive properties, including high scalability, high read/program speed, high energy efficiency, and programming voltages comparable with the power supply of typical neuromorphic chips \cite{kuzum2013Nanotech,ielmini2019Nanotech}. While non-volatile memristive devices have already shown promising results when used as non-volatile synapses, volatile memristive devices that can keep track of recent neural activity and implement synapses with biologically compatible time constants, are still largely unexplored. Indeed, volatile properties have so far mainly been investigated in reservoir computing applications \cite{du2017NatComm,midya2019AdvIntellSys,zhong2021NatComm}.

In this work, we use C\,/\,HfO$\mathrm{_{2}}$\,/\,Ag memristive devices featuring a high ON/OFF ratio of 10$^8$ and tunable time constants in the range of milliseconds to seconds \cite{wang2021AdvIntellSys,covi2021TED-I}, which are two desirable features for our application. The switching mechanism relies on the formation and spontaneous dissolution of a silver \gls{cf} across the switching layer. The relaxation, namely retention time, i.e., the time it takes to the \gls{cf} to dissolve, has been shown to depend on the diameter of the \gls{cf} itself: the thicker the filament, the longer the retention time \cite{wang2019TED-I,wang2019TED-II,wang2021TED-II}. The advantages of this technology are manifold: the retention times are electrically tunable and the time information is physically located in the \gls{cf}, thus consuming a negligible power and area on the chip. Moreover, the probability to switch the device on is also electrically tunable, thus adding a further degree of freedom in controlling the dynamics of the \gls{wm}. These features are exploited to demonstrate store and recall of patterns in a small scale \gls{wm} hardware, then the system is scaled up in simulations to demonstrate the use of volatile memristive devices in a large-scale biologically inspired model of synaptic \gls{wm} and in an associative symbolic \gls{wm}.


